\begin{textsample}{POS Dim 1 – vem\_ed\_28 – Score 54.00 – vem\_ed\_28/t147.txt}  \label{ex:f1_pos_004}
PCIs na colocação \textbf{profissional} dos tecnólogos Fatec Bragança Paulista A globalização transformou profundamente a educação \textbf{superior}, pois \textbf{tornou} a internacionalização uma \textbf{necessidade} para as instituições \textbf{que} buscam formar \textbf{profissionais} preparados para os desafios de um \textbf{mercado} de trabalho interconectado.
No \textbf{entanto}, a mobilidade física, muitas vezes associada ao \textbf{intercâmbio} tradicional, não \textbf{é} acessível para \textbf{todos} os estudantes, \textbf{seja} por \textbf{questões} financeiras, familiares ou outras limitações.
Nesse cenário, iniciativas como os Projetos Colaborativos Internacionais da Unidade do Ensino Superior de Graduação do Centro Paula Souza (PCIs / Cesu), o COIL (Collaborative Online International Learning) e o BRaVE (Brazilian Virtual Exchange), têm se mostrado alternativas valiosas, ao \textbf{permitirem} \textbf{que} alunos de diferentes \textbf{partes} do mundo colaborem em projetos acadêmicos sem a \textbf{necessidade} de se deslocarem fisicamente. \textbf{continuação} De \textbf{acordo} com Deardorff (2006), a internacionalização deve \textbf{ser} integrada ao currículo de maneira \textbf{que} desenvolva \textbf{competências} interculturais nos estudantes, \textbf{sendo} a \textbf{interação} com culturas diversas um \textbf{fator} \textbf{essencial} para o \textbf{desenvolvimento} dessas \textbf{habilidades}, \textbf{seja} de \textbf{forma} física ou virtual.
A autora \textbf{defende} \textbf{que}, embora o \textbf{intercâmbio} presencial \textbf{seja} frequentemente idealizado, as \textbf{interações} virtuais também proporcionam um enriquecimento acadêmico \textbf{significativo} ao \textbf{permitir} \textbf{que} os alunos compartilhem e adquiram novas \textbf{perspectivas} globais, o \textbf{que} contribui diretamente para seu crescimento intercultural. (Deardorff, 2006, p. 104).
Na Jornada de PCIs 2024, realizada em 6 de dezembro, nossa apresentação explorou o \textbf{impacto} dos Intercâmbios Virtuais, com \textbf{foco} \textbf{específico} na experiência dos alunos da Fatec Bragança Paulista, especialmente do curso de Análise e Desenvolvimento de Sistemas.
O diretor da unidade, professor Dr.
Vagner Donizeti Tavares Ferreira, \textbf{é} um grande incentivador dos projetos de Internacionalização em Casa.
Ele \textbf{reconhece} \textbf{que}, embora muitos dos estudantes enfrentem dificuldades financeiras \textbf{que} os impedem de participação em \textbf{programas} de mobilidade física, os Projetos Colaborativos Internacionais (PCIs / Cesu) \textbf{permitem} \textbf{que} participem de iniciativas internacionais de \textbf{forma} remota e adquiram \textbf{competências} globais e interculturais.
Os PCIs realizados pelos professores do Núcleo de Estudos da Linguagem (Nelf) da Fatec Bragança Paulista com universidades do Chile, por exemplo, fez com \textbf{que} os alunos brasileiros se \textbf{conectassem} com seus colegas chilenos por meio de atividades \textbf{colaborativas} e, no decorrer delas, praticassem três línguas: português, espanhol e inglês.
Esses projetos não apenas \textbf{ampliam} o \textbf{horizonte} acadêmico dos alunos, mas também os preparam para um \textbf{mercado} de trabalho globalizado.
A experiência de colaborar com estudantes de diferentes países contribui para o \textbf{desenvolvimento} de \textbf{habilidades} como comunicação intercultural, trabalho em equipe em \textbf{ambientes} diversificados e resolução de problemas em \textbf{contextos} globais. \textbf{continuação} Guth e Rubin (2022) discutem como o uso de plataformas de mobilidade virtual, como o COIL, facilita a aprendizagem \textbf{colaborativa} entre estudantes de diferentes países ao \textbf{criar} uma rede global de aprendizagem.
Eles argumentam \textbf{que} esses projetos \textbf{promovem} a \textbf{troca} de \textbf{conhecimentos} acadêmicos e ajudam os alunos a desenvolverem uma \textbf{compreensão} mais profunda das diferentes culturas e \textbf{realidades}, o \textbf{que} \textbf{é} \textbf{essencial} para a formação de \textbf{profissionais} globais.
Ao integrar essas práticas à formação \textbf{local}, os PCIs proporcionam uma educação mais inclusiva, acessível e \textbf{conectada} com as demandas do mundo moderno.
As equipes mistas compostas por alunos brasileiros e chilenos nos PCIs realizados na Fatec Bragança Paulista têm se \textbf{destacado} não apenas pelo \textbf{aprendizado} técnico, mas pela \textbf{capacidade} de interagir em um \textbf{ambiente} multicultural e multidisciplinar.
O uso de múltiplos \textbf{idiomas} no \textbf{processo} de comunicação, além de facilitar a \textbf{troca} de \textbf{conhecimentos}, também contribui para a melhoria das \textbf{competências} \textbf{linguísticas} dos alunos, \textbf{essenciais} para a atuação \textbf{profissional} em um cenário internacional.
Em suma, os PCIs oferecem uma oportunidade única de internacionalização para estudantes \textbf{que}, de outra \textbf{forma}, não poderiam participar de \textbf{intercâmbios} tradicionais.
Eles \textbf{ampliam} o acesso ao \textbf{conhecimento}, \textbf{fortalecem} a formação acadêmica e \textbf{tornam} os alunos mais competitivos no \textbf{mercado} de trabalho globalizado.
Ao \textbf{estabelecer} parcerias com instituições de ensino superior de outros países e \textbf{promover} \textbf{intercâmbios} virtuais, a Fatec de Bragança Paulista demonstra seu compromisso com uma educação internacionalizada, inclusiva e \textbf{alinhada} com as exigências do \textbf{futuro}.
Referências DEARDOFF, D.
K.
Identification and assessment of intercultural competence as a student outcome of internationalization.
Journal of Studies in International Education, v. 10, n.3, p. 241-266, Fall 2006.
Disponível em: https://doi.org/10.1177/1028315306287002.
Acesso em: 24 fev. 2025.
GUTH, S.; RUBIN, J.(eds.).
The Guide to COIL Virtual Exchange: Implementing, Growing, and Sustaining Collaborative Online International Learning.
New York: Routledge, 2022.
Disponível em: https://doi.org/10.4324/9781003447832.
Acesso em: 24 fev. 2025.

% matched lemmas: acordo, alinhar, ambiente, ampliar, aprendizado, capacidade, colaborativo, competência, compreensão, conectar, conhecimento, contexto, continuação, criar, defender, desenvolvimento, destacar, entanto, específico, essencial, estabelecer, fator, foco, forma, fortalecer, futuro, habilidade, horizonte, idioma, impacto, interação, intercâmbio, linguístico, local, mercado, necessidade, parte, permitir, perspectiva, processo, profissional, programa, promover, que, questão, realidade, reconhecer, ser, significativo, superior, todo, tornar, troca
\end{textsample}
